\documentclass[12pt,spanish,mexico]{article}
\usepackage[utf8]{inputenc}
\PassOptionsToPackage{hyphens}{url}
\usepackage{hyperref}
\usepackage[es-nodecimaldot]{babel} % nombre de opción corregido
\usepackage{amsthm, amssymb, amsmath, amsfonts, bbm, mathtools}
\usepackage{csquotes}
\usepackage{enumerate}
\usepackage[shortlabels]{enumitem}
\usepackage{titling}
\usepackage{blindtext}
\usepackage{lmodern} 
\usepackage{float}
\usepackage{mathrsfs}
\usepackage[capitalise]{cleveref}
\usepackage{caption} 
\usepackage{graphicx}
\usepackage{chngcntr}
\usepackage{enumitem}
\usepackage[ruled,vlined]{algorithm2e}

\counterwithin{figure}{section}
\numberwithin{equation}{section}
\usepackage[vmargin=1in, hmargin=0.5in]{geometry}  

\setlength{\topmargin}{-50pt}
\setlength{\headheight}{38.23112pt}
\setlength{\headsep}{24pt}

\makeatletter
\def\iftitle{\ifx\thetitle \empty \else }
\def\safetitle{\iftitle{\thetitle}\fi}
\def\headertitle{%
    \ifdefined\hw{\hw}\fi%
    \ifdefined\hw \iftitle{ -- }\fi%
    \safetitle%
}
\def\@maketitle{%
  \begin{center}%
Tarea 4 - Problemas inversos - 25/03/2026 \\
Muñoz Macías Christian Jaime 
    \vskip 1em%
    \hrule
  \end{center}%
  \par
  \vskip 1.5em%
}
\makeatother
\usepackage{amsmath}
\usepackage{tikz}

\newcommand\RADot{%
  \mathord{%
    \mspace{1mu}%
    \text{\radot}%
    \mspace{1mu}%
  }%
}

\newcommand\radot{%
    \tikz[line cap=round,x=1ex,y=1ex,line width=0.3pt]
    {\draw (0,1) |- (1,0) (0.55,0) arc(0:90:0.55); \fill (0.23,0.23) circle (0.05);}%
}

%% encabezado personalizado
\usepackage{fancyhdr}
%% estilo plain personalizado
\fancypagestyle{plain}{%
    \fancyhf{} % limpia encabezado y pie de página
    \renewcommand{\headrulewidth}{0pt}
    \setlength{\headheight}{0pt}
    \fancyfoot[R]{\thepage}
}
%% estilo para todas las páginas
\pagestyle{fancy}
\fancyhf{}
\fancyhead[C]{Maestría en Matemáticas Aplicadas}
\fancyfoot[R]{\thepage}

%%%%%%%%%%%%%%%%%%%%%%%%
%% DEFINICIONES ÚTILES %%
%%%%%%%%%%%%%%%%%%%%%%%%
%% definiciones matemáticas
\newcommand{\N}{\mathbb{N}}
\newcommand{\Z}{\mathbb{Z}}
\newcommand{\Ha}{\mathcal{H}}
\newcommand{\R}{\mathbb{R}}
\newcommand{\Q}{\mathbb{Q}}
\newcommand{\C}{\mathbb{C}}
\newcommand{\PP}{\mathbb{P}}

\newtheorem{proposition}{Proposición}
\newtheorem{theorem}{Teorema}
\newtheorem{lemma}{Lema}
\newtheorem{corollary}{Corolario}

\theoremstyle{definition}
\newtheorem{problem}{Problema}
\newtheorem{definition}{Definición}[section]
\newtheorem*{example}{Ejemplo}

\theoremstyle{remark}
\newtheorem*{observation}{Observación}
\newtheorem*{note}{Nota}

\newenvironment{solution}{\begin{proof}[Solución]}{\end{proof}}
\newenvironment{dem}{\begin{proof}[Demostración]}{\end{proof}}
\renewcommand{\qedsymbol}{$\blacksquare$}

%% comandos útiles
\newcommand{\eps}{\varepsilon}
\newcommand{\To}{\longrightarrow}
\newcommand{\Eu}{\varPhi}
\newcommand{\Mapsto}{\longmapsto}

%% delimitadores emparejados
\DeclarePairedDelimiter{\set}{\{}{\}}
\DeclarePairedDelimiter{\floor}\lfloor\rfloor
\DeclarePairedDelimiter{\ceil}\lceil\rceil
\DeclarePairedDelimiter{\abs}\lvert\rvert
\DeclarePairedDelimiter{\gen}\langle\rangle
\DeclareMathOperator{\sgn}{sgn}
\DeclareMathOperator{\ord}{ord}

%% operador
\newcommand{\op}[1]{\operatorname{#1}}

%% etiquetas del documento
\title{Tarea 4 - Problemas inversos - 25/03/2026}
\def\course{Materia}
\def\university{CIMAT - MsC. Maestría en matemáticas aplicadas}
\author{Muñoz Macías Christian Jaime}
\def\email{\url{christian.munoz@cimat.mx}}
\date{\today}

%% permitir saltos de línea en ecuaciones
\allowdisplaybreaks[4]

\begin{document}

\maketitle

\section*{Ecuación de Poisson}
Consideremos el problema de valores a la frontera para la ecuación de Poisson
\[
-u''(x) = f(x), \quad x \in (0, L), \quad u(0) = u(L) = 0. \tag{1}
\]
La solución se puede escribir como una integral usando la función de Green $G(x, \xi)$:
\[
u(x) = \int_{0}^{L} G(x, \xi) f(\xi) \, d\xi,
\]
donde la función de Green está dada por:
\[
G(x, \xi) = \begin{cases} 
\dfrac{x(L - \xi)}{L}, & \text{si } x \le \xi, \\[10pt]
\dfrac{\xi(L - x)}{L}, & \text{si } x > \xi. 
\end{cases}
\]
Entonces, la solución general del problema es:
\[
u(x) = \int_{0}^{x} \frac{\xi(L-x)}{L} f(\xi) \, d\xi + \int_{x}^{L} \frac{x(L-\xi)}{L} f(\xi) \, d\xi. \tag{2}
\]

\section*{Ejercicios}

\begin{enumerate}
\item Implementa la cuadratura de la regla del punto medio compuesta para discretizar el problema (2) suponiendo que $L = 1$.

\begin{solution}
Definimos el paso de discretización $h$ como 
\[
h = \frac{1}{n}
\]
donde $n$ es el número de subintervalos. La cuadratura de la regla del punto medio compuesta está definida como
\[
\xi_i^* = \left(i - \frac{1}{2}\right)h, \quad i = 1, \dots, n
\]
donde $\xi_i^*$ son los puntos de evaluación para $f(\xi)$, y la integral se aproxima por
\[
\int_{0}^{1} g(\xi)\,d\xi \approx h\sum_{i=1}^{n} g(\xi_i^*)
\]

Aplicando esto a la ecuación (2) con $L=1$:
\[
u(x) = \int_{0}^{x} \xi(1-x) f(\xi) \, d\xi + \int_{x}^{1} x(1-\xi) f(\xi) \, d\xi
\]

Discretizando:
\[
u(x) \approx h \sum_{i: \xi_i^* \le x} \xi_i^*(1-x) f(\xi_i^*) + h \sum_{i: \xi_i^* > x} x(1-\xi_i^*) f(\xi_i^*)
\]

Si consideramos un conjunto de puntos de observación $x_j$ para $j = 1, \dots, m$, entonces:
\[
u(x_j) \approx h \sum_{i=1}^{n} G(x_j, \xi_i^*) f(\xi_i^*)
\]
donde
\[
G(x_j, \xi_i^*) = \begin{cases}
\xi_i^*(1-x_j), & \text{si } x_j \le \xi_i^* \\
x_j(1-\xi_i^*), & \text{si } x_j > \xi_i^*
\end{cases}
\]

Esto se puede escribir como un sistema lineal:
\[
\mathbf{u} \approx A \mathbf{f}
\]
donde $\mathbf{u} \in \mathbb{R}^m$, $\mathbf{f} \in \mathbb{R}^n$, y la matriz $A \in \mathbb{R}^{m \times n}$ tiene elementos:
\[
\left[A\right]_{ji} = h \begin{cases}
\xi_i^*(1-x_j), & x_j \le \xi_i^* \\
x_j(1-\xi_i^*), & x_j > \xi_i^*
\end{cases}
\]
\end{solution}
\item Haz un modelo a priori de $f(x)$ mediante un Gaussian Markov Random Field usando el artículo de Bardsley como referencia. Considera que la condición de frontera de $f$ es conocida. Nota: Al construir la distribución a priori, es necesario añadir un offset (desplazamiento) a la media; de lo contrario, la media a priori será cero y el modelo no capturará adecuadamente la información disponible sobre $f$.

\begin{solution}
Se puede considerar que es un problema inverso lineal:
\[
b = Ax + \eta
\]
donde $b \in \mathbb{R}^m$ son los datos observados, $x \in \mathbb{R}^n$ es el vector desconocido, $A \in \mathbb{R}^{m \times n}$ es la matriz del modelo, y $\eta$ es ruido gaussiano.

La verosimilitud está dada por:
\[
p(b|x) \propto \exp\left(-\frac{\lambda}{2}\|Ax - b\|^2\right)
\]
donde $\lambda = 1/\sigma^2$ es la precisión del ruido.

Para el prior, consideramos un GMRF con media $\mu$:
\[
p(x) \propto \exp\left(-\frac{1}{2}(x-\mu)^T Q (x-\mu)\right)
\]
donde $Q$ es la matriz de precisión. En una malla uniforme con paso $h$, para un GMRF de primer orden:
\[
[Q]_{ij} = \frac{\delta}{h^2}
\begin{cases}
n_i, & i = j \\
-1, & j \in \partial_i \\
0, & \text{en otro caso}
\end{cases}
\]
donde $\partial_i$ son los vecinos del nodo $i$, $n_i = |\partial_i|$ es el número de vecinos, y $\delta > 0$ es un parámetro de regularización.

Para puntos interiores: $n_i = 2$, $Q_{ii} = 2\delta/h^2$, $Q_{i,i\pm1} = -\delta/h^2$.
Para puntos en la frontera: $n_i = 1$, $Q_{ii} = \delta/h^2$, y solo un vecino.

La distribución posterior es:
\[
p(x|b) \propto p(b|x)p(x) \propto \exp\left(-\frac{1}{2}(x-\mu_{\text{post}})^T \Sigma_{\text{post}}^{-1} (x-\mu_{\text{post}})\right)
\]
\end{solution}
\item Implementa la solución bayesiana del problema inverso de determinar
\[
f(x) = x
\]
dadas observaciones ruidosas de
\[
u(x) = \frac{1}{6}x(1-x)(1+x)
\]
con ruido aditivo gaussiano con relación de señal a ruido 100.

\item Calcula la media y la covarianza posterior de forma analítica.

\begin{solution}
Consideramos que 
\[
x|b \sim \mathcal{N}(\mu_{\text{post}}, \Sigma_{\text{post}})
\]
entonces $p(x|b)$ es
\[
p(x|b)=\frac{1}{(2\pi)^{n/2}|\Sigma_{\text{post}}|^{1/2}}\exp\left(-\frac{1}{2}(x-\mu_{\text{post}})^T \Sigma_{\text{post}}^{-1} (x-\mu_{\text{post}})\right)
\]
de su forma logarítmica
\[
\log p(x|b)=-\frac{n}{2}\log(2\pi)-\frac{1}{2}\log|\Sigma_{\text{post}}|-\frac{1}{2}(x-\mu_{\text{post}})^T \Sigma_{\text{post}}^{-1} (x-\mu_{\text{post}})
\]
si conocemos la proporción de $\log p(x|b)$
\begin{align*}
\log p(x|b) &\propto \log p(b|x)+ \log p(x) \\
&\propto -\frac{\lambda}{2}\|Ax-b\|^2 -\frac{n}{2} \log(2\pi)+ \frac{1}{2}\log|Q| -\frac{1}{2}(x-\mu)^T Q (x-\mu)
\end{align*}
haciendo un reacomodo algebraico
\begin{align*}
\log p(x|b) &\propto -\frac{n}{2} \log(2\pi)+ \frac{1}{2}\log|Q|-\frac{\lambda}{2} \langle Ax-b, Ax-b\rangle-\frac{1}{2} \langle x-\mu,Q (x-\mu)\rangle
\end{align*}
la parte de interés es 
\begin{align*}
&=-\frac{\lambda}{2} \langle Ax-b, Ax-b\rangle-\frac{1}{2} \langle x-\mu,Q (x-\mu)\rangle \\
&=-\frac{\lambda}{2} \left(\langle Ax, Ax\rangle-\langle Ax, b\rangle-\langle b, Ax\rangle +\langle b, b\rangle\right)-\frac{1}{2}\left(\langle x,Q (x-\mu)\rangle-\langle \mu,Q (x-\mu)\rangle \right)\\
&=-\frac{\lambda}{2} \left( x^T A^T A x - x^T A^T b - b^T A x + b^T b \right)-\frac{1}{2}\left(x^T Q x - x^T Q \mu - \mu^T Q x +\mu^T Q \mu \right)
\end{align*}
como $Q$ es simétrica entonces $- x^T Q \mu - \mu^T Q x = -2 x^T Q \mu$
\begin{align*}
&= -\frac{1}{2} x^T (\lambda A^T A + Q) x + (\lambda b^T A + \mu^T Q) x - \frac{\lambda}{2}b^T b - \frac{1}{2}\mu^T Q \mu
\end{align*}
Esta expresión es una forma cuadrática en $x$:
\[
\log p(x|b) \propto -\frac{1}{2} x^T (\lambda A^T A + Q) x + (\lambda b^T A + \mu^T Q) x + \text{constante}
\]
Comparando con la forma cuadrática de una distribución normal:
\[
-\frac{1}{2} (x-\mu_{\text{post}})^T \Sigma_{\text{post}}^{-1} (x-\mu_{\text{post}}) = -\frac{1}{2} x^T \Sigma_{\text{post}}^{-1} x + \mu_{\text{post}}^T \Sigma_{\text{post}}^{-1} x + \text{constante}
\]
Identificamos:
\[
\Sigma_{\text{post}}^{-1} = \lambda A^T A + Q
\]
\[
\mu_{\text{post}}^T \Sigma_{\text{post}}^{-1} = \lambda b^T A + \mu^T Q
\]
Por lo tanto:
\[
\Sigma_{\text{post}} = (\lambda A^T A + Q)^{-1}
\]
\[
\mu_{\text{post}} = \Sigma_{\text{post}} (\lambda A^T b + Q \mu)
\]

\item Muestra tu distribución posterior y compara la media posterior numérica tomando en cuenta el tiempo integrado de autocorrelación.
\begin{figure}[H]
    \centering
        \includegraphics[width=0.7\linewidth]{C:/Users/chris/OneDrive/Escritorio/python/CIMAT/problemas inversos/Tarea_4/T4_g1.pdf}
        \caption{Resultados punto 3,4,5}
    \end{figure}
\end{solution}

\item Calcula la constante de normalización de manera analítica, y numéricamente de la manera siguiente:


\begin{enumerate}
\item[(a)] Usa la función \texttt{scipy.stats.gaussian\_kde} para ajustar una distribución gaussiana $G(f)$ a las muestras que obtuviste.

\item[(b)] Aproxima
\[
\begin{aligned}
z_0 &= \int \pi(u|f)\pi(f)\,df \\
&= \int \frac{\pi(u|f)\pi(f)}{G(f)} G(f)\,df \\
&\approx \frac{1}{N} \sum_{j=1}^{N} \frac{\pi(u|f_j)\pi(f_j)}{G(f_j)}
\end{aligned}
\]
donde $\{f_j\}_{j=1}^N$ son muestras de $G(f)$.
\end{enumerate}
\begin{solution}
Media posterior (primeros 5 valores): [1.29075943 2.35771351 0.87590171 0.71068738 2.4701287 ]
Covarianza posterior (primeros 5x5):
[[ 0.02743117 -0.03298458 -0.01009924 -0.02295709  0.01102738]
 [-0.03298458  0.05561902 -0.00505292  0.02789583 -0.02264902]
 [-0.01009924 -0.00505292  0.02372022  0.00668544  0.00670482]
 [-0.02295709  0.02789583  0.00668544  0.02727986 -0.01940718]
 [ 0.01102738 -0.02264902  0.00670482 -0.01940718  0.02348941]]
log(z0) analitico: -370.799642
log(z0) numerico: -1599.520470
diferencia: 1228.720828
\end{solution}

\end{enumerate}


\end{document}